\documentclass[11pt]{article}
\usepackage{amsmath,amssymb,amsthm,hyperref}
\newtheorem{theorem}{Theorem}
\newtheorem{proposition}{Proposition}
\title{What the Exponential Pole Count Does and Does Not Supply}
\author{Joint Cell Head Compression}
\date{24 November 2026}
\begin{document}
\maketitle

This note concerns the candidate in \texttt{joint-cell-filtered-kernels.tex}.
It does not prove a new asymptotic bound for the original head complexity.
The equal-step identities and scalar word budget below are kernel checked
in \texttt{Results/JointCellEqualStepCapacity.lean}. The moment argument is
a written proof, not a Lean declaration. Unequal steps, several independent
tokens, and affine row numerators remain outside the equal-step conclusion.

For positive weights $w_1,\ldots,w_c$, put $s_w=\sum_iw_i$ and
\[
 K_w(d)=\sum_{T\subseteq[c]}\frac{(-1)^{c-|T|}}{d-\sum_{i\in T}w_i},
 \qquad d>s_w.
\]

\begin{theorem}[Equal-step collapse]
If $w_i=w$ for every $i$, then
\[
 K_w(d)=\sum_{j=0}^c\frac{(-1)^{c+j}\binom cj}{d-jw}.
\]
Consequently, suppose each of $a$ scalar output bits is a bias plus a sum
of $q$ equal-step filters, with arbitrary bit-dependent amplitudes, shifts,
and step sizes. If all denominators are positive at ordered tokens and the
tokens yield $m+1$ distinct strict sign words, then
\[
 m\leq aq(c+1).
\]
In particular, decoding every $a$-bit word requires
$2^a-1\leq aq(c+1)$.
\end{theorem}

\begin{proof}
Group subsets by cardinality. Each bit has a partial-fraction expansion
with at most $q(c+1)$ poles, counting multiplicities. Clearing its positive
denominator gives a polynomial of degree at most $q(c+1)$. Consecutive
distinct words differ in at least one coordinate. The corresponding
polynomial has a root in that token gap. The gaps are disjoint, so the sum
of polynomial degrees is at least their number. The Lean declarations are
\texttt{jointCell\_filtered\_kernel\_equal\_steps},
\texttt{jointCell\_bitwise\_cauchy\_word\_budget}, and
\texttt{jointCell\_equal\_step\_filter\_word\_budget} in namespace
\texttt{HeadComplexity.AxisZeroGluing}.
\end{proof}

For fixed $q$ and $c\leq n$, the word length is only $O(\log n)$.
More generally the efficiency $a/q$ is $O(\log(c+2))$, even when $q$ grows.
Here is an elementary verification of the latter statement. Let
$L=\lceil\log_2(c+2)\rceil$. If $a\leq4L$, the assertion follows from
$q\geq1$. Otherwise use
\[
 \frac aq\leq\frac{a^2(c+1)}{2^a-1}
 \leq\frac{2a^2(c+1)}{2^a}
 \leq\frac{32L^2}{8^L}\leq4L.
\]
The penultimate inequality uses $c+1\leq2^L$ and the decrease of
$u^2/2^u$ for integers $u\geq4$; the last uses $8L\leq8^L$.
Thus increasing the number of equal-step filters cannot rescue the
proposed efficiency. This conclusion concerns a shared scalar token.
It is not a lower bound for independent multivariate token decoders.

\begin{proposition}[Fixed unequal weights cannot hide arbitrarily many
leading moments at infinity]
Let $w,v\in(0,\infty)^c$ and define
\[
 G_w(z)=\frac{z^{c+1}}{c!\prod_iw_i}K_w(z+s_w/2),
 \qquad z>s_w/2.
\]
If $G_w(z)-G_v(z)=O(z^{-2c-2})$ as $z\to+\infty$, then $w$ and $v$
have the same multiset of coordinates and $G_w=G_v$. If their multisets
are different, their first unequal Laurent coefficient has order at most
$2c$.
\end{proposition}

\begin{proof}
Iterating the fundamental theorem of calculus in the mixed difference
gives, for independent $U_i$ uniform on $[-w_i/2,w_i/2]$ and
$X_w=\sum_iU_i$,
\[
 \frac{K_w(z+s_w/2)}{c!\prod_iw_i}
   =\mathbb E(z-X_w)^{-c-1}.
\]
The expansion is convergent for $z>s_w/2$, and symmetry gives
\[
 G_w(z)=\sum_{r\geq0}\binom{c+2r}{2r}
              \mathbb E X_w^{2r}\,z^{-2r}.
\]
Equality through order $2c$ therefore equates the first $c$ even moments.
Moments determine cumulants triangularly. The moment generating function
of $X_w$ is
\[
 \mathbb E e^{tX_w}=\prod_i\frac{\sinh(w_it/2)}{w_it/2},
 \qquad
 \kappa_{2r}(X_w)=\frac{B_{2r}}{2r}\sum_iw_i^{2r}.
\]
The Bernoulli numbers $B_{2r}$ are nonzero. Their generating function and
nonvanishing sign are recorded in the NIST DLMF,
\href{https://dlmf.nist.gov/24.2.E1}{24.2.1} and
\href{https://dlmf.nist.gov/24.2.E2}{24.2.2}.
Thus the power sums of the $c$ numbers $w_i^2$ and $v_i^2$ agree through
degree $c$. Newton's identities give equality of their monic root
polynomials. Positivity then gives equality of the original multisets.
\end{proof}

The proposition only treats a fixed pair of normalized centered kernels
at infinity. It does not bound zeros near a finite center, allow weights
to vary with the compactness parameter, or handle affine numerators.
It therefore does not close the unequal-step route. Numerical searches
have found ratios with several turns, but have not established an
exponential number of usable sign intervals or a uniform upper bound.
Moreover, a usable codebook must fit a token interval of width $\delta$
while retaining a denominator margin larger than the sum of all column
and bit slope changes. One-hot positivity alone is insufficient.

\end{document}
